\documentclass[11pt]{article}
\usepackage[margin=1.05in]{geometry}
\usepackage{amsmath,amssymb}
\usepackage{xcolor}
\usepackage[colorlinks=true,linkcolor=violet,urlcolor=teal]{hyperref}
\usepackage{enumitem}
\newcommand{\entry}[2]{\medskip\noindent\textbf{#1}\quad #2\par}
\title{\bfseries Annotated bibliography for the almost-all non-halting theorem}
\author{BB(6) research program}
\date{2026-09-08}
\begin{document}
\maketitle

\noindent This bibliography accompanies the literature review and the
manuscript. Each entry states the reference, the precise result, and its role
in the proof. Bibliographic details were verified against primary sources
(publisher pages, NUMDAM, AMS, IMPAN, Mathnet, and in two cases the article
PDFs directly); residual uncertainties are flagged in \emph{italics}. Nothing
here is cited for an effect its source does not actually provide --- in
particular we distinguish what supplies the ACIM from what supplies the
\emph{rate}, and what is finite-measure from what is infinite-measure.

\section{Foundations and the Collatz template}

\entry{Tao (2022).}{\emph{Almost all orbits of the Collatz map attain almost
bounded values}, Forum of Mathematics, Pi \textbf{10}, e12 (arXiv:1909.03562).
The theorem: for any $f\to\infty$, $\mathrm{Col}_{\min}(N)<f(N)$ for almost all
$N$ in \emph{logarithmic} density. The method passes to the Syracuse (odd-to-odd)
map, models the $2$-adic valuations by a ``Syracuse random variable,'' and
proves a local limit theorem controlled by a bound on the characteristic
function of a skew random walk on $\mathbb Z/3^n\mathbb Z$. \textbf{Role:} the
template for the entire program --- reduce the arithmetic map, model the
increments, prove a local-limit/equidistribution estimate, conclude an
almost-all statement. We transcribe its architecture (statistical model $+$
quantitative rate; first-passage decomposition) while replacing its $3$-adic
harmonic analysis by a suspension/Dolgopyat argument, since our cost is
multiplicative rather than a Birkhoff sum on a compact group.}

\entry{Korec (1994); Krasikov--Lagarias (2003).}{\emph{A density estimate for
the $3x+1$ problem}, Math.\ Slovaca \textbf{44}, 85--89; and \emph{Bounds for
the $3x+1$ problem using difference inequalities}, Acta Arith.\ \textbf{109},
237--258 (arXiv:math/0205002). The natural-density predecessors: Korec bounds
the eventual minimum below $N^{\theta}$, $\theta>\log3/\log4$, for almost all
$N$; Krasikov--Lagarias show $\ge x^{0.84}$ of integers below $x$ provably reach
$1$. \textbf{Role:} they calibrate ``almost all'' --- these are power-law
partial results, and Tao's advance (which we follow) was to reach density one
via a genuine LLT rather than a sieve. \emph{Applegate--Lagarias (Math.\ Comp.\
64, 1995) and Lagarias's annotated $3x{+}1$ bibliography are companion pointers;
exact pages/arXiv ids there are medium-confidence.}}

\entry{Lagarias--Weiss (1992).}{\emph{The $3x+1$ problem: two stochastic
models}, Ann.\ Appl.\ Probab.\ \textbf{2}, 229--261. The rigorous stochastic
heuristics that the density results approximate. \textbf{Role:} the ancestor of
the ``biased random walk'' intuition that our theorem makes rigorous for the
renewal map. \emph{Pages high-confidence; confirm before print.}}

\entry{Aaronson (2020).}{\emph{The Busy Beaver Frontier}, ACM SIGACT News
\textbf{51}(3), 2020 (also ECCC TR20-115). A survey of $\mathrm{BB}(n)$: growth,
independence from ZFC of $\mathrm{BB}(n)$ for explicit $n$, and small machines
encoding open conjectures. \textbf{Role:} frames the meta-mathematical stakes
and the bridge from $\mathrm{BB}(6)$ machines to Collatz-type dynamics.
\emph{Pages 32--54 medium-confidence.}}

\entry{The bbchallenge Collaboration (2024--2025).}{The Coq-verified proof that
$\mathrm{BB}(5)=47{,}176{,}870$ (Coq-BB5, led by the contributor ``mxdys,''
announced July 2024; companion paper arXiv:2509.12337), and \emph{Antihydra},
the six-state machine whose halting is equivalent to a Collatz-type first-passage
event for $n\mapsto\lfloor3n/2\rfloor$ from $n=8$. \textbf{Role:} the concrete
$\mathrm{BB}(6)$ context --- Antihydra is the canonical exhibit of a machine
whose single-orbit behaviour is Collatz-hard while its statistics are the
natural object of an almost-all theorem. \emph{Authorship is collective/
pseudonymous; the halting-probability figure is a heuristic, not a theorem; the
current $\mathrm{BB}(6)$ champion value should be quoted from the bbchallenge
wiki if used.} \emph{As of 2025, Tao (2022) stands essentially alone as a
rigorous almost-all-log-density result for a Collatz-type map; no credible
extension to another arithmetic system was found --- so a port to a
$\mathrm{BB}(6)$/Antihydra-derived map would be the first.}}

\section{Invariant measures, mixing, and the transfer operator}

\entry{Lasota--Yorke (1973).}{\emph{On the existence of invariant measures for
piecewise monotonic transformations}, Trans.\ AMS \textbf{186}, 481--488. The
Lasota--Yorke inequality $\mathrm{Var}(\mathcal Lf)\le\alpha\,\mathrm{Var}(f)
+C\|f\|_1$, $\alpha<1$, giving an ACIM with BV density. \textbf{Role:}
the ancestor of every LY estimate we use.}

\entry{Rychlik (1983).}{\emph{Bounded variation and invariant measures}, Studia
Math.\ \textbf{76}, 69--80. Extends Lasota--Yorke to \emph{countably many}
branches under summable variation: unique ACIM, BV density, quasi-compact
transfer operator. \textbf{Role:} this is what makes our countable-branch map's
ACIM and untwisted spectral gap free --- the accumulation point $\rho^\ast$ is a
point of strong expansion, so no infinite-measure machinery is needed.
\emph{Attribution caution: Rychlik supplies the countable-branch ACIM and BV
framework; the exponential decay \emph{rate} is classically co-attributed to
Hofbauer--Keller (below), which we cite for the rate.}}

\entry{Hofbauer--Keller (1982).}{\emph{Ergodic properties of invariant measures
for piecewise monotonic transformations}, Math.\ Z.\ \textbf{180}, 119--140.
Exponential decay of correlations for piecewise-expanding interval maps.
\textbf{Role:} the mixing-\emph{rate} citation companion to Rychlik.
\emph{Pages high-confidence; not DOI-verified in this pass.}}

\entry{Keller--Liverani (1999).}{\emph{Stability of the spectrum for transfer
operators}, Ann.\ Scuola Norm.\ Sup.\ Pisa (4) \textbf{28}, 141--152. Isolated
spectral data are stable under a perturbation small in a weak norm (need not be
small in the strong BV norm). \textbf{Role:} the theorem that transfers the
finite-branch gap to the countable map via truncation, and that absorbs the
$O(1/a)$ model-to-exact perturbation of the \emph{bulk} branches. (Not the tool
for the native-resolution scale perturbation --- see the manuscript's item-1
discussion --- which is handled instead by applying the gap per exact-scale
map.)}

\entry{Baladi (2018).}{\emph{Dynamical Zeta Functions and Dynamical Determinants
for Hyperbolic Maps: A Functional Approach}, Ergebnisse der Mathematik 3.\ Folge
\textbf{68}, Springer. The standard modern account of transfer-operator spectral
theory via anisotropic spaces. \textbf{Role:} background reference for the
functional-analytic framework. \emph{Note: a \emph{maps} reference; it does not
contain the Dolgopyat oscillatory-cancellation estimate, which is separate.}}

\entry{Mayer (1990).}{\emph{On the thermodynamic formalism for the Gauss map},
Comm.\ Math.\ Phys.\ \textbf{130}, 311--333. Transfer-operator spectral theory
of the continued-fraction Gauss map. \textbf{Role:} the Gauss-type exemplar our
induced map is modelled on. \emph{If a continued-fraction/modular equidistribution
statement were wanted instead, Series (JLMS 1985) is the alternative; the plain
Gauss--Kuzmin--Wirsing spectral picture (Wirsing, Acta Arith.\ 24, 1974) is the
lighter option.}}

\section{The Dolgopyat method: uniform spectral gaps}

\entry{Dolgopyat (1998).}{\emph{On decay of correlations in Anosov flows},
Annals of Math.\ (2) \textbf{147}, 357--390 (verified from the article PDF).
Introduces the ``Dolgopyat method'': an $L^2$ cancellation estimate on twisted
transfer operators along the imaginary frequency, driven by uniform
non-integrability (UNI) of the temporal-distance function, giving exponential
mixing. \textbf{Role:} the origin of the uniform large-frequency gap we need.}

\entry{Dolgopyat (1998), Prevalence of rapid mixing.}{Ergodic Theory Dynam.\
Systems \textbf{18}(5), 1097--1114. States UNI in reusable \emph{symbolic} form
as a checkable condition on the roof cocycle. \textbf{Role:} the precise
citation for ``UNI as a hypothesis one verifies,'' as opposed to the Annals
paper's flow-specific use.}

\entry{Naud (2005).}{\emph{Expanding maps on Cantor sets and analytic
continuation of zeta functions}, Ann.\ Sci.\ ENS (4) \textbf{38}, 116--153
(verified from the article PDF). \textbf{The best-matched reference.} For
uniformly expanding Markov maps with real-analytic roof, the non-local
integrability condition NLI (Def.\ 2.1) is that the temporal-distance function
$\varphi_{\xi,\eta}$ have non-vanishing derivative somewhere --- in the analytic
case simply $\varphi_{\xi,\eta}\not\equiv0$, i.e.\ non-proportional branch roofs.
Theorem~2.3: under NLI, $\|\mathcal L_s^n\|_{C^1}\le C|\mathrm{Im}\,s|^{1+\epsilon}
\rho^n$, $\rho<1$, \emph{uniformly in $\mathrm{Im}\,s$}; Theorem~1.7: a zero-free
strip for the zeta function. \textbf{Role:} our UNI verification --- branch forms
$L_k(\xi)=(r_k-p_k)+(s_k-q_k)\xi$ pairwise non-proportional (slopes $-\rho_k$
distinct) --- is exactly Naud's $\varphi_{\xi,\eta}\not\equiv0$ on our arithmetic,
so Theorem~2.3 is the gap we invoke.}

\entry{Stoyanov (2011).}{\emph{Spectra of Ruelle transfer operators for Axiom A
flows}, Nonlinearity \textbf{24}, 1089--1120 (arXiv:0810.1126). The
higher-dimensional Axiom-A-flow analogue of Naud's estimate under a UNI/Diophantine
temporal-distance condition. \textbf{Role:} corroborating generality.
\emph{Several closely-titled Stoyanov papers exist; pin by arXiv id.}}

\entry{Avila--Gou\"ezel--Yoccoz (2006).}{\emph{Exponential mixing for the
Teichm\"uller flow}, Publ.\ Math.\ IH\'ES \textbf{104}, 143--211. Exponential
mixing via a Dolgopyat estimate over a \emph{countable-alphabet} (Rauzy--Veech)
roof. \textbf{Role:} the methodological precedent that the Dolgopyat estimate
survives countably many branches --- moral support for our \S(infinite) reduction.}

\entry{Ruelle (1976); Pollicott (1985); Parry--Pollicott (1990).}{\emph{Zeta-functions
for expanding maps and Anosov flows}, Invent.\ Math.\ \textbf{34}, 231--242;
\emph{On the rate of mixing of Axiom A flows}, Invent.\ Math.\ \textbf{81},
413--426; \emph{Zeta functions and the periodic orbit structure of hyperbolic
dynamics}, Ast\'erisque \textbf{187--188}. The transfer-operator/zeta framework,
Pollicott--Ruelle resonances, and the RPF/non-lattice theory. \textbf{Role:}
the framework citations; Naud invokes Parry--Pollicott Thm 6.3 for the
non-lattice continuation. \emph{Pollicott 1985 pages secondary-verified.}}

\section{Local limit theorems}

\entry{Nagaev (1957); Guivarc'h--Hardy (1988).}{\emph{Some limit theorems for
stationary Markov chains}, Theory Probab.\ Appl.\ \textbf{2}, 378--406; and
\emph{Th\'eor\`emes limites pour une classe de cha\^\i nes de Markov et
applications aux diff\'eomorphismes d'Anosov}, Ann.\ IHP Probab.\ Statist.\
\textbf{24}(1), 73--98. The origin of the perturbed-operator (eigenvalue)
approach to the CLT, and its adaptation to hyperbolic dynamics with LLT and
renewal theorems. \textbf{Role:} the ``origin'' citations for the method that
turns our spectral gap into a local limit theorem.}

\entry{Hennion--Herv\'e (2001).}{\emph{Limit Theorems for Markov Chains and
Stochastic Properties of Dynamical Systems by Quasi-Compactness}, Springer LNM
\textbf{1766}. The definitive treatment of the Nagaev--Guivarc'h method:
CLT/LLT/Berry--Esseen from quasi-compactness plus analytic perturbation of the
leading eigenvalue $\lambda(i\xi)$. \textbf{Role:} the machinery citation
bridging our BV/Lipschitz gap to the LLT output.}

\entry{Aaronson--Denker (2001).}{\emph{Local limit theorems for partial sums of
stationary sequences generated by Gibbs--Markov maps}, Stoch.\ Dyn.\ \textbf{1}(2),
193--237. LLT for Birkhoff sums of Lipschitz observables over Gibbs--Markov
(countable-partition, big-image, bounded-distortion) maps, under aperiodicity
(non-lattice) and finite variance. \textbf{Role:} the off-the-shelf LLT for our
log-growth cocycle over the induced Gauss-type map. \emph{Cite this long
published title, not the short-title preprint.}}

\section{Countable-branch thermodynamic formalism (context, not load-bearing)}

\entry{Sarig (1999).}{\emph{Thermodynamic formalism for countable Markov shifts},
Ergodic Theory Dynam.\ Systems \textbf{19}(6), 1565--1593. Gurevich pressure,
variational principle, and RPF theory for countable-state shifts under the BIP
condition. \textbf{Role:} the symbolic backdrop for coding our countable-branch
map; not load-bearing given Rychlik.}

\entry{Aaronson--Denker--Urba\'nski (1993).}{\emph{Ergodic theory for Markov
fibred systems and parabolic rational maps}, Trans.\ AMS \textbf{337}(2),
495--548. Thermodynamic formalism for countable Markov fibred systems.
\textbf{Role:} structural context for the fibred/countable-branch setting.}

\entry{Sarig (2002); Gou\"ezel (2004); Melbourne--Terhesiu (2012).}{\emph{Subexponential
decay of correlations}, Invent.\ Math.\ \textbf{150}, 629--653; \emph{Sharp
polynomial estimates for the decay of correlations}, Israel J.\ Math.\ \textbf{139},
29--65; \emph{Operator renewal theory and mixing rates for dynamical systems with
infinite measure}, Invent.\ Math.\ \textbf{189}, 61--110. Operator-renewal theory
for slow (polynomial/subexponential) mixing --- Sarig/Gou\"ezel finite-measure,
Melbourne--Terhesiu infinite-measure. \textbf{Role:} cited only to say they are
\emph{not} needed: our tail is geometric and our measure finite, so the
truncation is elementary. Included so a reader is not left wondering whether the
countable branches force slow-mixing machinery. \emph{They do not.}}

\section*{On confidence}
Two load-bearing items (Dolgopyat 1998; Naud 2005) were verified from the
article PDFs, including the theorem statements used. All other primary entries
have venue/volume/year confirmed against publisher or archive pages; italicized
flags mark page ranges or identifiers confirmed only from consistent secondary
metadata, and the handful of companion references not separately DOI-checked.
The two anomalous ``post-Tao extension'' preprints surfaced during the search
were judged unverifiable and are deliberately \emph{excluded}.

\end{document}
